\chapter{Center and scale in shrinker stability}\label{sec:shrinker-stability}
Allowing the Gaussian center and scale to vary changes the stability question. We keep this finite-dimensional adjustment explicit. The source is \cite{CM12}, Theorem 4.14, Lemma 4.23, and Sections 5--6 and 11.1. Write $[u]=\int_\Sigma u(4\pi)^{-n/2}e^{-|x|^2/4}\,dA$ and use $L$ from Definition \ref{def:shrinker-jacobi}.

\begin{proposition}[First variation with a moving Gaussian center and scale]\label{prop:gaussian-joint-first}
Let $\Sigma_s$ be a smooth compactly supported variation of a smooth two-sided properly immersed boundaryless hypersurface with polynomial area growth. Let the center and scale satisfy $c_0=0$, $t_0=1$, $c'_0=y$, $t'_0=h$, and let the initial normal speed be $f$. For
\[
 F_{c,t}(\Sigma)=(4\pi t)^{-n/2}\int_\Sigma e^{-|x-c|^2/(4t)}\,dA
\]
we have
\[
 \left.\frac{d}{ds}\right|_0 F_{c_s,t_s}(\Sigma_s)
 =[f\Phi+h\alpha+\beta],\qquad
 \alpha=\tfrac{|x|^2}{4}-\tfrac n2,\quad \beta=\tfrac{\langle x,y\rangle}{2}.
\]
If $\Sigma$ is also a complete connected self-shrinker, this derivative is zero.
\end{proposition}
\begin{proof}
Differentiate the area factor and the logarithm of the Gaussian density. The shape terms give $f\Phi$ by Proposition \ref{prop:gaussian-compact-variation}; differentiating the scale gives $h\alpha$, and differentiating the center gives $\beta$. Outside the changing compact set, Gaussian moments from Lemma \ref{lem:gaussian-polynomial-moments} dominate the parameter derivatives uniformly near $s=0$. At a shrinker, $\Phi=0$, $[\alpha]=0$ and $[\beta]=0$ by Proposition \ref{prop:shrinker-basic-moments}. This is the normalized first-variation formula underlying \cite{CM12}, Sections 3--4.
\end{proof}

\begin{theorem}[The simultaneous shape, center, and scale Hessian]\label{thm:shrinker-joint-hessian}
Let $\Sigma^n\subset\mathbb R^{n+1}$ be a connected complete smooth properly immersed two-sided self-shrinker without boundary and with polynomial area growth. For a smooth compactly supported shape variation with initial normal speed $f$, and center/scale derivatives $y\in\mathbb R^{n+1}$ and $h\in\mathbb R$ at $(c,t)=(0,1)$,
\[
 \left.\frac{d^2}{ds^2}\right|_0 F_{c_s,t_s}(\Sigma_s)
 =Q(f)+2h[fH]-h^2[H^2]+[f a_y]-\tfrac12[a_y^2],
 \qquad a_y=\langle y,\nu\rangle.
\]
Here $[u]=\int_\Sigma u(4\pi)^{-n/2}e^{-|x|^2/4}\,dA$, $Q$ is Definition \ref{def:shrinker-jacobi}, and $h=t'_0$ is the derivative of the Gaussian time scale. This is \cite{CM12}, Theorem 4.14, pp.~777--778.
\end{theorem}
\begin{proof}
At a critical point the second derivative depends only on initial velocities: differences in accelerations are paired with the vanishing first variation. We may therefore use linear center and scale paths and a normal parametrization to compute. At $s=0$, differentiating the general versions of the three factors in Proposition \ref{prop:gaussian-joint-first} gives
\[
 \Phi'=-Lf+hH+\tfrac12a_y,\quad
 \alpha'=fH-\beta-\tfrac h2(|x|^2-n),\quad
 \beta'=\tfrac12fa_y-\tfrac12|y|^2-h\beta.
\]
For example $\Phi_s=H_s-\langle x_s-c_s,\nu_s\rangle/(2t_s)$ and $\alpha_s=|x_s-c_s|^2/(4t_s^2)-n/(2t_s)$. At the shrinker the logarithmic derivative of the weighted area element is $h\alpha+\beta$. The product rule consequently yields
\[
 F''=\left[-fLf+2hfH+fa_y-\tfrac{h^2}{2}(|x|^2-n)
 -2h\beta-\tfrac12|y|^2+(h\alpha+\beta)^2\right].
\]
Now $[\beta]=[\alpha\beta]=0$ by Proposition \ref{prop:shrinker-basic-moments}. Proposition \ref{prop:shrinker-fourth-moment} gives $[\alpha^2]=n[1]/2-[H^2]$, while $[|x|^2]=2n[1]$. Thus the quadratic scale terms reduce to $-h^2[H^2]$. Finally $[\beta^2]=[|y^T|^2]/2$, again by the moment identities, so the center terms reduce to $-[a_y^2]/2$. All differentiations outside compact shape support involve only polynomial Gaussian moments, already integrable by the stated growth hypothesis.
\end{proof}

\begin{definition}[F-stability and its quantifiers]\label{def:shrinker-f-stability}
A shrinker satisfying the hypotheses of Theorem \ref{thm:shrinker-joint-hessian} is $F$-stable if, for every $f\in C_c^\infty(\Sigma)$, there exist a center derivative $y$ and a scale derivative $h$ for which that second variation is nonnegative. It is $F$-unstable if some such $f$ makes the second variation negative for every $y,h$. This is \cite{CM12}, Section 4.3. Fixing $y=h=0$ instead gives the different question in Theorem \ref{thm:shrinker-fixed-instability}. Entropy uses the supremum over all centers and positive scales, Definition \ref{def:mcf-entropy}; local $F$-stability does not by definition settle that global supremum problem. Theorem \ref{thm:shrinker-entropy-lowering} supplies the entropy-decrease implication for an $F$-unstable shrinker without a Euclidean line factor.
\end{definition}

\begin{proposition}[Removing geometric modes on a closed shrinker]\label{prop:shrinker-optimized-hessian}
Let $\Sigma$ be a connected smooth closed two-sided immersed self-shrinker. Set $\mathcal T=\{\langle y,\nu\rangle:y\in\mathbb R^{n+1}\}$. Every smooth $f$ has a unique orthogonal decomposition
\[
 f=aH+b+g,\qquad b\in\mathcal T,\quad g\perp H,\quad g\perp\mathcal T
\]
in Gaussian $L^2$. Its joint Hessian is
\[
 Q(g)-(a-h)^2[H^2]-\tfrac12[(b-a_y)^2].
\]
Hence its maximum over $h,y$ is $Q(g)$, attained when $h=a$ and $a_y=b$.
\end{proposition}
\begin{proof}
First $H$ cannot vanish identically: if it did, $\Delta|x|^2=2n$ on the closed manifold, contradicting its integral being zero. Corollary \ref{cor:shrinker-mode-orthogonality} makes $\mathbb RH$ orthogonal to the finite-dimensional space $\mathcal T$, giving the decomposition. Symmetry of $L$ and its equations $LH=H$, $Lb=b/2$ give
\[
 Q(f)=-a^2[H^2]-\tfrac12[b^2]+Q(g),\quad
 [fH]=a[H^2],\quad [fa_y]=[ba_y].
\]
Insert these in Theorem \ref{thm:shrinker-joint-hessian} and complete the two squares. By the definition of $\mathcal T$, the required vector $y$ exists; uniqueness of $y$ is unnecessary. No spectral assumption about a noncompact surface has entered.
\end{proof}

\begin{proposition}[The round sphere is F-stable]\label{prop:shrinker-sphere-f-stable}
The sphere $S^n(\sqrt{2n})\subset\mathbb R^{n+1}$, $n\geq1$, is $F$-stable. For the remainder $g$ orthogonal to constants and translation modes,
\[
 Q(g)\geq\frac1n[g^2].
\]
The estimate uses the classical spherical harmonic spectrum as an external spectral input, as in \cite{CM12}, Lemma 4.23.
\end{proposition}
\begin{proof}
Here $x^T=0$, $|A|^2=1/2$, $H=\sqrt{n/2}$, and $L=\Delta+1$. The external spherical harmonic decomposition gives eigenvalues $k(k+n-1)/(2n)$ of $-\Delta$, $k=0,1,\ldots$. Constants are the $k=0$ space; the $k=1$ space consists of restrictions of linear functions, precisely the normal translation modes up to a constant factor. On their orthogonal complement the eigenvalues are at least $(n+1)/n$. Expansion in the orthogonal spherical harmonics therefore gives $[|\nabla g|^2]-[g^2]\geq[g^2]/n$. Proposition \ref{prop:shrinker-optimized-hessian} supplies center and scale adjustments making the full Hessian nonnegative. The proof does not contradict fixed-center instability.
\end{proof}

\begin{theorem}[Sign-changing mean curvature gives closed F-instability]\label{thm:shrinker-changing-sign-instability}
If a connected smooth closed two-sided immersed self-shrinker has scalar mean curvature $H$ that changes sign, it is $F$-unstable. The argument uses compact elliptic spectral theory as an external input; it is the instability step in \cite{CM12}, Section 6, pp.~782--783.
\end{theorem}
\begin{proof}
The external principal-eigenfunction theorem for a smooth self-adjoint elliptic operator on a connected closed manifold says that the largest eigenvalue $\lambda_0$ of $L$ is simple and admits a strictly positive smooth eigenfunction $u$. This is Corollary 5.15 of \cite{CM12}, written with $Lu=\lambda_0u$ rather than $Lu=-\mu_1u$. Since $LH=H$ and $H$ changes sign, $H$ cannot be the principal eigenfunction; simplicity gives $\lambda_0>1$. Weighted symmetry then makes $u$ orthogonal to $H$ and every translation mode, whose eigenvalue is $1/2$. The joint Hessian for $f=u$ is
\[
 -\lambda_0[u^2]-h^2[H^2]-\tfrac12[a_y^2]<0
\]
for every $h,y$. Compactness makes $u$ an admissible compactly supported variation. This proves instability, without using or claiming the separate classification of mean-convex shrinkers.
\end{proof}

\begin{theorem}[Nontrivial round cylinders are F-unstable]\label{thm:shrinker-cylinder-f-unstable}
For $1\leq k<n$, the self-shrinker $S^k(\sqrt{2k})\times\mathbb R^{n-k}$ is $F$-unstable for compactly supported normal variations. This is the cylinder case of \cite{CM12}, Section 11.1, pp.~816--817. In particular the roles of cylinders in entropy and singularity theory cannot be summarized as $F$-stability under Definition \ref{def:shrinker-f-stability}.
\end{theorem}
\begin{proof}
Write $z\in\mathbb R^{n-k}$ for the axial coordinate. The cylinder has $H=\sqrt{k/2}$, $|A|^2=1/2$, and $Lz_1=z_1/2$. Choose smooth radial cutoffs $0\leq\eta_R(z)\leq1$, equal to one for $|z|\leq R$, zero for $|z|\geq2R$, and with $|\nabla\eta_R|\leq C/R$. Set $f_R=\eta_Rz_1$. This has compact support because the spherical factor is compact. Its oddness in $z_1$ gives $[f_R]=0$ exactly. The normal depends only on the spherical factor, and its integral on that sphere is zero, so $[f_Ra_y]=0$ for every $y$.

The Gaussian product measure has all polynomial moments. Therefore $f_R\to z_1$ and $\nabla f_R\to\nabla z_1$ in Gaussian $L^2$: the cutoff errors are supported in Gaussian tails and $|z_1\nabla\eta_R|\leq C|z|/R$. Since $|A|^2+1/2=1$,
\[
 Q(f_R)\longrightarrow [|\nabla z_1|^2-z_1^2]
 =-\tfrac12[z_1^2]<0.
\]
The last equality follows from the one-dimensional Gaussian identity $[z_1^2]=2[1]$, obtained by integrating the derivative of $z_1e^{-z_1^2/4}$ on the real line. Choose one large $R$ with $Q(f_R)<0$. Its joint Hessian is $Q(f_R)-h^2[H^2]-[a_y^2]/2$, strictly negative for every center and scale derivative. Exact oddness makes this conclusion uniform in $h$, without passing a nonuniform limit through an optimization.
\end{proof}
